% Generated draft: requires mathematical and source review.
\documentclass[11pt]{article}
\usepackage[margin=25mm]{geometry}
\usepackage{fontspec}
\setmainfont{Latin Modern Roman}
\usepackage{amsmath,amssymb,mathtools,unicode-math}
\setmathfont{Latin Modern Math}
\usepackage{enumitem,needspace,xurl,hyperref}
\setlength{\parindent}{0pt}
\setlength{\parskip}{5pt}
\setlength{\emergencystretch}{3em}
\newcommand{\N}{\mathbb{N}}
\newcommand{\Z}{\mathbb{Z}}
\newcommand{\Q}{\mathbb{Q}}
\newcommand{\R}{\mathbb{R}}
\newcommand{\C}{\mathbb{C}}
\newcommand{\F}{\mathbb{F}}
\newcommand{\A}{\mathbb{A}}
\newcommand{\PP}{\mathbb{P}}
\newcommand{\E}{\mathbb{E}}
\newcommand{\eps}{\varepsilon}
\newcommand{\dd}{\,\mathrm{d}}
\begin{document}
\section*{UnsolvedMath ID: 30000275}
\textbf{Source problem code:} catalogue entry
\subsection{Definitions and mathematical statement}
Let \(k\) be a number field, \(\A_k\) its adele ring, and \(C\) the curve over \(k\) intended in the source. View \(C(k)\) diagonally in \(C(\A_k)\). The original statement takes \(X\) to be a residue class in the adelic points disjoint from \(C(k)\), whereas the cleaned statement writes \(X\in C(\A_k)\), the notation for an individual adelic point. These are different input types and are not silently conflated.

For a subset interpretation, let \(X\subseteq C(\A_k)\) be an admissible residue class with \(X\cap C(k)=\varnothing\). Given an unramified covering \(f:D\to C\), let \(f_j:D_j\to C\) range over all its twists defined over \(k\) in the intended twisting theory. The weak question asks whether one can choose \(f\) such that
\[
 \forall j,\qquad f_j\bigl(D_j(\A_k)\bigr)\cap X=\varnothing,
\]
or equivalently \(X\cap\bigcup_j f_j(D_j(\A_k))=\varnothing\). The same one covering must exclude the whole prescribed \(X\) for every twist; choosing a different covering for each twist is not the assertion. The stronger question asks whether that covering can always be chosen abelian. Under the cleaned point interpretation, the corresponding condition is \(X\notin f_j(D_j(\A_k))\) for every \(j\), with \(X\notin C(k)\); the original residue-class hypothesis is not then automatically recovered.

In the usual finite-descent setup a covering is finite \(\acute{e}\)tale, and twists are specified by a finite group scheme or torsor structure. For a torsor under \(G\), the standard twist labels lie in \(H^1(k,G)\), with appropriately twisted structure groups and maps. This explains one possible precise setup only: the source does not give a torsor structure, finiteness, connectedness or geometric connectedness conditions, the group used for twists, or whether ``abelian'' concerns a geometric deck group or a group scheme over \(k\). The meaning and topology of an adelic residue class, including its places and moduli, are likewise unspecified. Smoothness, projectivity and any genus restrictions on the curve are absent.

The 2026-08-21 literature assessment calls these weak and strong covering conjectures, reports important special cases, and says they would imply that the Brauer--Manin obstruction is the only obstruction for curves. This implication is retained as an attributed source assessment; its exact hypotheses and formulation are not supplied and are not proved here. The open label and the verification note's editorial recasting remain dated metadata.
\subsection{Short English statement}
Can one unramified cover exclude a rational-point-free adelic residue class from the adelic images of every twist, and can such a cover be abelian? The cleaned wording instead treats the residue class as a point. Both readings and their different exclusion conditions are recorded, with the covering and twisting conventions left for source confirmation.
\subsection{Sources}
\begin{itemize}
\item Local UnsolvedMath registry, permanent ID 30000275, source code catalogue entry, ``Abelian Covering Conjectures for Curves over Number Fields,'' including the original residue-class wording, cleaned point notation and both covering questions.
\item \emph{Explicit Methods in Number Theory (2006)}, Oberwolfach Reports, pp.\ 1799--1866, DOI 10.4171/owr/2005/32: \url{https://doi.org/10.4171/owr/2005/32}. The citation uses 2006, while the DOI and proposed-year metadata use 2005.
\item Supplied reference ``Stoll finite descent obstructions'': \url{https://doi.org/10.2140/ant.2007.1.349}.
\item Literature triage checked 2026-08-21; local record updated 2026-09-16. No external reference has been visited or independently checked here.
\end{itemize}
\end{document}
